\subsection{Integral jets and why abelian freeness is insufficient}
For a finite free $\Z$-algebra $B$, put
$C_\pm(\mathcal D_B)=\mathcal D_q(B,a)/(c\mp1)\mathcal D_q(B,a)$.
The scalar Tate calculation extends to this coefficient ring.
\begin{proposition}\label{irjet:thm:finite-free-koszul}
With $\overline B=B/2B$ and the reduced active sequence
$\overline{\boldsymbol\delta}$, the two integer-torsion submodules are
\[
 \operatorname{Tor}_{\Z}C_+(\mathcal D_B)
 \simeq\bigoplus_{j\text{ odd}} H_j(K_{\overline B}(\overline{\boldsymbol\delta})),
 \qquad
 \operatorname{Tor}_{\Z}C_-(\mathcal D_B)
 \simeq\bigoplus_{j\text{ even}} H_j(K_{\overline B}(\overline{\boldsymbol\delta})).
\]
These are isomorphisms of $B$-modules. Every torsion class is killed by two.
For $q>2$ and $T=\operatorname{rank}_{\Z}B$, the free abelian rank of
both quotients is $T\varphi(q)/2$.
\end{proposition}
\begin{proof}
The underlying distribution lattice is free over $\Z$ by the all-ring
basis theorem. Thus the torsion lemma still applies. In the fixed-point
periodic complex, multiplication by two is injective on $B$ and its
free chain modules. Its adjacent two-term blocks resolve reduction modulo
two, exactly as in \eqref{ir:eq:integer-tate-koszul}; the bounded chain
width gives the same parity sums as $B$-modules. Over $B\otimes\Q$,
reflection eigenspaces are exact direct summands. The fixed-point Euler
count in the balanced-rank proof is multiplied by $T$ and remains zero,
giving the asserted free ranks. This argument does not assume that the
reflected torsion-free part is a free $B$-module.
\end{proof}

\subsection{An exact Smith formula for arbitrary finite integral jets}
\label{irjet:sec:jets}
Let
\[
 B_L=\Z[u]/(u^L),\quad L\ge1,\qquad
 a_p(u)\in B_L,
\]
and reduce coefficients modulo two. In
$\overline B_L=\mathbb F_2[u]/(u^L)$, form the parameters
\eqref{ir:eq:active-koszul}. For an element $f$ of this ring, define
\[
 \nu(f)=\min\{j:[u^j]f\ne0\},\qquad \nu(0)=L,
 \quad v=\min_{f\in\boldsymbol f}\nu(f).
\]
Here $\boldsymbol f=\overline{\boldsymbol\delta}$ and $r(q)$ is the active-prime count in \eqref{ir:eq:r-h}. For $q>2$, the list is nonempty.

\begin{theorem}[Minimum-valuation jet formula]
\label{irjet:thm:jets}
For $q>2$, the two reflected jet quotients have the following Smith
decomposition as abelian groups:
\begin{equation}
 C_\pm(\mathcal D_{B_L})\simeq
 \Z^{L\varphi(q)/2}\oplus(\Z/2\Z)^{v\,2^{r(q)-1}}.
 \label{irjet:eq:jet-Smith}
\end{equation}
More precisely, their integer-torsion submodules are each isomorphic,
as $B_L$-modules, to
\begin{equation}
 \bigl(B_L/(2,u^v)\bigr)^{2^{r(q)-1}}.
 \label{irjet:eq:jet-torsion-module}
\end{equation}
The direct-sum decomposition in \eqref{irjet:eq:jet-Smith} is a statement over
$\Z$, not a claim that the torsion-free part is free over $B_L$.
\end{theorem}
\begin{proof}
In the principal ideal ring $\overline B_L$, choose a parameter of least
valuation. It is a unit times $u^v$ unless every parameter is zero,
which is the same argument with $v=L$. A change of basis in the free
module defining the Koszul complex transforms the parameter list into
$(u^v,0,\ldots,0)$: first divide by its unit factor, then subtract its
multiples from the other entries. All these are invertible operations
over $\overline B_L$.

The complex is now the tensor product of the two-term complex given by
multiplication by $u^v$ and $r(q)-1$ zero-differential exterior factors.
Both the kernel and the cokernel of multiplication by $u^v$ are
isomorphic to $\overline B_L/(u^v)$; the kernel is the ideal
$u^{L-v}\overline B_L$. Consequently
\[
 H_k\bigl(K_{\overline B_L}(\boldsymbol f)\bigr)
 \simeq\bigl(\overline B_L/(u^v)\bigr)^{\binom{r(q)}{k}}.
\]
Summing either degree parity in Theorem~\ref{irjet:thm:finite-free-koszul} proves
\eqref{irjet:eq:jet-torsion-module}. Each copy has $\mathbb F_2$-dimension $v$,
so Lemma~\ref{ir:lem:torsion-tate} gives the torsion multiplicity in
\eqref{irjet:eq:jet-Smith}.

The unreflected underlying lattice has rank $L\varphi(q)$.
The alternating fixed-point trace used in
Theorem~\ref{ir:thm:integer-reflection} is multiplied by $L$, and still
vanishes. Thus each rational eigenspace has dimension
$L\varphi(q)/2$. This proves the free rank and completes the Smith
calculation.
\end{proof}

\begin{example}[Five jet coefficients at level fifteen]
Take $L=5$, $a_3(u)=1+u^2$, and $a_5(u)=1+2u+u^3$.
The binary parameters are $u^2,u^3$, so $v=2$ and $r(q)=2$. Therefore
\[
 C_\pm(\mathcal D_{B_5})\simeq\Z^{20}\oplus(\Z/2\Z)^4.
\]
The constant jets $a_3=a_5=1$ would instead give ten two-torsion factors.
The constant terms of the weights agree; their higher binary
coefficients change the integral reflection structure.
\end{example}

\begin{example}[The two-adic jet parameter matters]
At $q=12$, $L=4$, let $a_2(u)=1+u$ and $a_3(u)=1+u^3$.
The two parameters are $u^3$ and $u+u^2$, not simply
$1-a_3$ and $1-a_2$. Thus $v=1$, $r(q)=2$, and
\[
 C_\pm(\mathcal D_{B_4})\simeq\Z^8\oplus(\Z/2\Z)^2.
\]
Although the parity of a \emph{scalar} two-prime weight does not affect
Theorem~\ref{ir:thm:integer-reflection}, a nonconstant two-prime jet can affect
Theorem~\ref{irjet:thm:jets}.
\end{example}

\subsubsection{Why the free abelian part need not be jet-free}
Take $q=3$, $L=2$, and $a_3=1+u$. In the unreflected basis
$e=e_0$, $f=e_{2/3}$, reflection satisfies
\[
 ce=e,\qquad cf=ue-f.
\]
Hence
\begin{equation}
 C_+(\mathcal D_{B_2})=B_2e\oplus B_2f\big/\langle ue-2f\rangle.
 \label{irjet:eq:nonflat}
\end{equation}
As an abelian group this has free generators $e,f$ and a torsion generator
$uf$ of order two. In its torsion-free quotient, multiplication by $u$
sends $e$ to $2f$ and $f$ to zero. This is not a free rank-one
$B_2$-module: the image of $u$ in a free rank-one module is a primitive
rank-one sublattice, whereas here the image is $2\Z f$.
Thus a claim of $B_L$-freeness for the reflected torsion-free part would
be false. By contrast, $C_-=B_2f\oplus(B_2/(2,u))e$ in this example.
Both signs have the same abelian Smith factors, but different module
extensions over the jet ring.

\subsubsection{Multivariate jets and arithmetic weights}
For a finite free $\Z$-algebra $B$ of rank $L$, the Koszul theorem remains
an exact answer. If $q>2$ and
$h_k=\dim_{\mathbb F_2} H_k(K_{B/2B}(\boldsymbol f))$, then
\[
 \operatorname{Tor}_\Z C_+\simeq(\Z/2\Z)^{\sum_{k\text{ odd}}h_k},\qquad
 \operatorname{Tor}_\Z C_-\simeq(\Z/2\Z)^{\sum_{k\text{ even}}h_k}.
\]
The two sums are equal: the Euler characteristic of the Koszul complex
is $L(1-1)^{r(q)}=0$. A closed minimum-valuation description is special
to the one-variable principal ideal ring; it must not be transferred
unchanged to a multivariate Artinian ring.

For actual positive integral polylogarithm orders, $A_p=p^{1-n}$ usually
belongs to $\Z[1/q]$, rather than $\Z$. The corresponding coefficient
ring must be stated. If $q$ is even, inverting $q$ also inverts two and
removes the reflected two-torsion. For odd $q$, every $p^{1-n}$ is one
modulo two, and the ordinary maximal multiplicity persists over
$\Z[1/q]$. These are statements about the chosen formal coefficient
lattice, not about independence of the evaluated periods.
